\scanpage{12}

% archive sheet number: 11

Catégorie $\unsure{\mathcal{C}|\mathscr{C}}$ avec unité à droite $e$, et produits

$G$ groupe dans $\mathcal{C}$ \qquad $\unsure{p|\rho}_G:G\times G\to G$

$E$ $\mathcal{C}$-espace : $\mathcal{C}$-groupe $G$ opérant \unsure{à|:} gauche \qquad $q_E:G\times E\to E$

On dit que $E$ est formellement hom. principal. Si l'hom

\[
G\times E\xrightarrow{(\deleted{\unsure{g}}q_E,p_2)}E\times E
\]

\margin{$x=g(x,y)y$}

\margin{\unsure{\S|$S$}}

est un \uline{isomorphisme}. Cela signifie aussi que l'on peut trouver un hom

\[
E\times E\xrightarrow{g_E}G\quad\unsure{\prime\prime|,}
\]

tel que

\diagram{\begin{tikzcd}[ampersand replacement=\&]
\text{(i)}\quad E\times E\arrow[r,"{(g,p_2)}"]\arrow[rr,bend right=20,"p_1"'] \& G\times E\arrow[r,"q_E"]\& E\quad\text{est commutatif}
\end{tikzcd}}

\diagram{\begin{tikzcd}[ampersand replacement=\&]
\text{(ii)}\quad G\times E\arrow[r,"{(q_E,p_2)}"]\arrow[rr,bend right=20,"{p_{\deleted{\unsure{2}}1}}"'] \& E\times E\arrow[r,"g_E"]\& G\quad\text{est commutatif}
\end{tikzcd}}

Si $\mathcal{C}$ est la catégorie des $A$-algèbres (commutatives avec unité), alors $G,E$ correspondent à des $A$-algèbres \deleted{\unsure{comm.}} $B,C$, et on a

\[
\left\{\begin{aligned}
p_{\deleted{\unsure{G}}B}&:B\longrightarrow B\otimes B\\
q_C&:C\longrightarrow B\otimes C
\end{aligned}\right.\tag*{(*)}
\]

et le fait qu'on ait \uline{formellement} un fibré hom. principal signifie que

\[
\varphi:C\otimes C\xrightarrow{(q_C,p_2)}B\otimes C
\]

est un isomorphisme, i.e. qu'il existe

\[
g_C:B\longrightarrow C\otimes C
\]
